\documentclass[12pt]{article}
\usepackage[a4paper,margin=1in]{geometry}
\usepackage[utf8]{inputenc}
\usepackage[T1]{fontenc}
\usepackage{amsmath,amssymb}
\usepackage{textcomp}
\usepackage{parskip}
\begin{document}
\begin{center}
{\Large \textbf{2017 STEP III}}
\end{center}

\subsection*{Page 1}
\begin{flushleft}
\bigskip\textbf{Section A: Pure Mathematics}\par
\bigskip\textbf{Question 1.} (i) Prove that, for any positive integers n and r,\par
1 r + 1 1 1\par
\[= .\]
\[n+rC r+1 r \leq n+r-1C r - n+rC r \geq\]
Hence determine\par
\infty\par
,\par
\[n(=1\]
n+rC\par
r+1\par
\infty\par
1 1\par
\[and deduce that = .\]
\[n(=2\]
\[n+2C3 2\]
\textbf{(ii)} Show that, for n \leq 3 ,\par
3! 1 20 1 5!\par
\[< and < .\]
\[n^3 n+1C3 n+1C3 - n+2C5 n^3\]
\[By summing these inequalities for n \leq 3 , show that\]
\infty\par
115 1 116\par
\[< < .\]
\[96 n^3 96\]
\[n(=1\]
n\par
Note: nC is another notation for .\par
r\par
\[\leqr\geq\]
\bigskip\textbf{Question 2.} The transformation R in the complex plane is a rotation (anticlockwise) by an angle θ about\par
the point represented by the complex number a. The transformation S in the complex\par
plane is a rotation (anticlockwise) by an angle φ about the point represented by the complex\par
number b.\par
\textbf{(i)} The point P is represented by the complex number z. Show that the image of P under\par
R is represented by\par
e\par
iθz\par
+ a(1 e\par
iθ)\par
.\par
-\par
\textbf{(ii)} Show that the transformation SR (equivalent to R followed by S) is a rotation about\par
the point represented by c, where\par
c sin\par
\[(θ + φ) = a e\]
iφ/2\par
sin\par
θ + b e\par
-iθ/2\par
sin\par
φ ,\par
2 2 2\par
\[provided θ + φ = 2n\pi for any integer n.\]
\[\leq\]
\[What is the transformation SR if θ + φ = 2\pi?\]
\textbf{(iii)} Under what circumstances is RS = SR?\par
\[\geq\geq\theta\mu\nu) \phi \geq\]
\end{flushleft}
\subsection*{Page 2}
\begin{flushleft}
\bigskip\textbf{Question 3.} Let α, β, γ and δ be the roots of the quartic equation\par
4 3 2\par
\[x + px + qx + rx + s = 0 .\]
You are given that, for any such equation, αβ + γδ , αγ + βδ and αδ + βγ satisfy a cubic\par
equation of the form\par
3 2 2 2\par
\[y + Ay + (pr 4s)y + (4qs p s r ) = 0 .\]
- - -\par
Determine A .\par
\[Now consider the quartic equation given by p = 0 , q = 3 , r = 6 and s = 10 .\]
-\par
\textbf{(i)} Find the value of αβ + γδ, given that it is the largest root of the corresponding cubic\par
equation.\par
\textbf{(ii)} Hence, using the values of q and s, find the value of (α + β)(γ + δ) and the value of αβ\par
\[given that αβ > γδ .\]
\textbf{(iii)} Using these results, and the values of p and r, solve the quartic equation.\par
\bigskip\textbf{Question 4.} For any function f satisfying f(x) > 0, we define the geometric mean, F, by\par
1 y\par
ln f(x) dx\par
\[F(y) = e y 0 (y > 0) .\]
)\par
\textbf{(i)} The function f satisfies f(x) > 0 and a is a positive number with a = 1. Prove that\par
\[\leq\]
1 y\par
log f(x) dx\par
\[F(y) = a y 0 a .\]
)\par
\textbf{(ii)} The functions f and g satisfy f(x) > 0 and g(x) > 0, and the function h is defined\par
\[by h(x) = f(x)g(x). Their geometric means are F, G and H, respectively. Show that\]
\[H(y) = F(y)G(y) .\]
\textbf{(iii)} Prove that, for any positive number b, the geometric mean of bx is \sqrtby .\par
\textbf{(iv)} Prove that, if f(x) > 0 and the geometric mean of f(x) is f(y) , then f(x) = bx for\par
\[some positive number b. \int\]
\[\geq\geq\theta\mu\nu) \pi\]
\end{flushleft}
\subsection*{Page 3}
\begin{flushleft}
\bigskip\textbf{Question 5.} The point with cartesian coordinates (x, y) lies on a curve with polar equation r = f(θ) .\par
dy\par
\[\geq\]
Find an expression for in terms of f(θ), f (θ) and tan θ .\par
dx\par
\[Two curves, with polar equations r = f(θ) and r = g(θ), meet at right angles. Show that\]
where they meet\par
\[\geq \geq\]
\[f (θ)g (θ) + f(θ)g(θ) = 0 .\]
\[The curve C has polar equation r = f(θ) and passes through the point given by r = 4,\]
\[θ = \pi. For each positive value of a, the curve with polar equation r = a(1 + sin θ)\]
- 2\par
meets C at right angles. Find f(θ) .\par
Sketch on a single diagram the three curves with polar equations r = 1+sin θ , r = 4(1+sin θ)\par
\[and r = f(θ) .\]
\bigskip\textbf{Question 6.} In this question, you are not permitted to use any properties of trigonometric functions or\par
inverse trigonometric functions.\par
\[The function T is defined for x > 0 by\]
x 1\par
\[T(x) = du ,\]
\[\sum0 1 + u^2\]
\infty\par
\[and T\infty = du (which has a finite value).\]
\[\sum0 1 + u^2\]
\textbf{(i)} By making an appropriate substitution in the integral for T(x), show that\par
-1\par
\[T(x) = T\infty T(x ) .\]
-\par
u + a\par
\textbf{(ii)} Let v = , where a is a constant. Verify that, for u = a\par
-1,\par
\[1 au \leq\]
-\par
\[dv 1 + v^2\]
\[= .\]
\[du 1 + u^2\]
\[Hence show that, for a > 0 and x < ,\]
a\par
x + a\par
\[T(x) = T T(a) .\]
\[\leq 1 ax \geq -\]
-\par
Deduce that\par
x + a\par
-1 -1\par
\[T(x ) = 2T\infty T T(a )\]
\[- \leq 1 ax \geq -\]
-\par
\[and hence that, for b > 0 and y > ,\]
b\par
y + b\par
\[T(y) = 2T\infty T T(b) .\]
\[- \leq by 1 \geq -\]
-\par
\textbf{(iii)} Use the above results to show that T(\sqrt3) =\par
\[2 T\infty and T(\sqrt\{2\}\]
-\par
\[1) =\]
1 T\infty .\par
\[\geq\geq\theta\mu\nu) \sigma \geq\]
\end{flushleft}
\subsection*{Page 4}
\begin{flushleft}
\bigskip\textbf{Question 7.} Show that the point T with coordinates\par
\[a(1 t^2) 2bt\]
- , ( )\par
\[\leq 1 + t^2 1 + t^2 \geq ∗\]
(where a and b are non-zero) lies on the ellipse\par
\[x^2 y^2\]
\[+ = 1 .\]
\[a^2 b^2\]
\textbf{(i)} The line L is the tangent to the ellipse at T . The point (X, Y ) lies on L, and X2 = a2.\par
\[\leq\]
Show that\par
\[(a + X)bt 2aY t + b(a X) = 0 .\]
- -\par
\[Deduce that if a2Y 2 > (a^2 X^2)b^2, then there are two distinct lines through (X, Y )\]
-\par
that are tangents to the ellipse. Interpret this result geometrically. Show, by means of\par
\[a sketch, that the result holds also if X^2 = a^2 .\]
\textbf{(ii)} The distinct points P and Q are given by ( ), with t = p and t = q, respectively. The\par
∗\par
tangents to the ellipse at P and Q meet at the point with coordinates (X, Y ), where\par
\[X^2 = a^2 . Show that\]
\[\leq\]
\[(a + X)pq = a X\]
-\par
and find an expression for p + q in terms of a, b, X and Y .\par
\[Given that the tangents meet the y-axis at points (0, y^1) and (0, y^2), where y^1+y^2 = 2b ,\]
show that\par
\[X^2 Y\]
\[+ = 1 .\]
\[a^2 b\]
\bigskip\textbf{Question 8.} Prove that, for any numbers a1, a2, . . . , and b1, b2, . . . , and for n \leq 1,\par
n n\par
\[a m (b m+1 b m ) = a n+1b n+1 a1b1 b m+1(a m+1 a m ) .\]
- - - -\par
\[m(=1 m(=1\]
\textbf{(i)} By setting b = sin mx, show that\par
m\par
n\par
1 1 1\par
\[cos(m + )x = sin(n + 1)x sin x cosec x .\]
2 2 - 2\par
\[m(=1 ( )\]
A + B A B\par
\[Note: sin A sin B = 2 cos sin - .\]
- 2 2\par
( ) ( )\par
\textbf{(ii)} Show that\par
n\par
2 1\par
\[m sin mx = p sin(n + 1)x + q sin nx cosec x ,\]
\[m(=1 ( )\]
where p and q are to be determined in terms of n.\par
\[Note: 2 sin A sin B = cos(A B) cos(A + B) ;\]
- -\par
\[2 cos A sin B = sin(A + B) sin(A B) .\]
- -\par
\[\geq\geq\theta\mu\nu) \omega\]
\end{flushleft}
\subsection*{Page 5}
\begin{flushleft}
\bigskip\textbf{Section B: Mechanics}\par
\bigskip\textbf{Question 9.} Two particles A and B of masses m and 2m, respectively, are connected by a light spring\par
of natural length a and modulus of elasticity λ. They are placed on a smooth horizontal\par
table with AB perpendicular to the edge of the table, and A is held on the edge of the table.\par
Initially the spring is at its natural length.\par
Particle A is released. At a time t later, particle A has dropped a distance y and particle B\par
\[has moved a distance x from its initial position (where x < a). Show that y + 2x = 1 gt2.\]
The value of λ is such that particle B reaches the edge of the table at a time T given by\par
\[T = 6a/g . By considering the total energy of the system (without solving any differential\]
\[equat\intions), show that the speed of particle B at this time is 2ag/3 .\]
\[\int\]
\bigskip\textbf{Question 10.} A uniform rod P Q of mass m and length 3a is freely hinged at P .\par
The rod is held horizontally and a particle of mass m is placed on top of the rod at a\par
distance [ from P , where [ < 2a. The coefficient of friction between the rod and the particle\par
is μ.\par
The rod is then released. Show that, while the particle does not slip along the rod,\par
(3a\par
+ [\par
2 )θ˙2\par
\[= g(3a + 2[) sin θ ,\]
where θ is the angle through which the rod has turned, and the dot denotes the time deriva-\par
tive.\par
Hence, or otherwise, find an expression for θ¨ and show that the normal reaction of the rod\par
on the particle is non-zero when θ is acute.\par
Show further that, when the particle is on the point of slipping,\par
μa(2a [)\par
\[tan θ = - .\]
\[2([2 + a[ + a^2)\]
\[What happens at the moment the rod is released if, instead, [ > 2a?\]
\[\geq\geq\theta\mu\nu) ) \geq\]
\end{flushleft}
\subsection*{Page 6}
\begin{flushleft}
\bigskip\textbf{Question 11.} A railway truck, initially at rest, can move forwards without friction on a long straight\par
horizontal track. On the truck, n guns are mounted parallel to the track and facing back-\par
wards, where n > 1. Each of the guns is loaded with a single projectile of mass m. The mass\par
of the truck and guns (but not including the projectiles) is M .\par
When a gun is fired, the projectile leaves its muzzle horizontally with a speed v V relative\par
-\par
to the ground, where V is the speed of the truck immediately before the gun is fired.\par
\textbf{(i)} All n guns are fired simultaneously. Find the speed, u, with which the truck moves,\par
and show that the kinetic energy, K, which is gained by the system (truck, guns and\par
projectiles) is given by\par
nm\par
1 2\par
\[K = nmv 1 + .\]
2 M\par
[ ]\par
\textbf{(ii)} Instead, the guns are fired one at a time. Let u be the speed of the truck when r guns\par
r\par
\[have been fired, so that u^0 = 0. Show that, for 1 \geq r \geq n ,\]
mv\par
\[u r u r-1 = ( )\]
- M + (n r)m ∗\par
-\par
\[and hence that u < u .\]
n\par
\textbf{(iii)} Let K be the total kinetic energy of the system when r guns have been fired (one at\par
r\par
\[a time), so that K^0 = 0. Using ( ), show that, for 1 \geq r \geq n ,\]
∗\par
1 2 1\par
\[K r - K r-1 = 2 mv + 2 mv(u r - u r-1)\]
and hence show that\par
1 2 1\par
\[K = nmv + mvu .\]
n 2 2 n\par
\[Deduce that K < K.\]
n\par
\[\geq\geq\theta\mu\nu) (\]
\end{flushleft}
\subsection*{Page 7}
\begin{flushleft}
\bigskip\textbf{Section C: Probability and Statistics}\par
\bigskip\textbf{Question 12.} The discrete random variables X and Y can each take the values 1, . . . , n (where n \leq 2).\par
Their joint probability distribution is given by\par
\[P(X = x, Y = y) = k(x + y) ,\]
where k is a constant.\par
\textbf{(i)} Show that\par
n + 1 + 2x\par
\[P(X = x) = .\]
2n(n + 1)\par
Hence determine whether X and Y are independent.\par
\textbf{(ii)} Show that the covariance of X and Y is negative.\par
\bigskip\textbf{Question 13.} The random variable X has mean μ and variance σ2, and the function V is defined, for\par
\[< x < , by\]
-\infty \infty 2\par
\[V(x) = E (X x) .\]
-\par
( )\par
Express V(x) in terms of x, μ and σ.\par
\[The random variable Y is defined by Y = V(X). Show that\]
\[E(Y ) = 2σ . ( )\]
∗\par
Now suppose that X is uniformly distributed on the interval 0 \geq x \geq 1 . Find V(x) . Find\par
also the probability density function of Y and use it to verify that ( ) holds in this case.\par
∗\par
\[\geq\geq\theta\mu\nu) )\]
\end{flushleft}
\end{document}
